\chapter{Programmi su più file e make}
\label{chap:make}

Finora ogni programma stava in un solo file. I programmi veri sono fatti di
molti file, a volte migliaia, e compilarli tutti a ogni modifica sarebbe
lentissimo. Questo capitolo mostra come si divide un programma in più file,
come si compilano separatamente e si collegano (riprendendo il
capitolo~\ref{chap:dal-problema}), e come il programma \code{make} automatizza
il lavoro, ricompilando solo ciò che serve.

\section{Perché più file}

Dividere un programma in più file sorgente ha diversi vantaggi:

\begin{itemize}
  \item \textbf{organizzazione}: ogni file raccoglie le funzioni che
        riguardano un aspetto del problema (l'input, i calcoli, l'output),
        come i capitoli di un libro;
  \item \textbf{riuso}: un file con funzioni utili (per esempio di
        geometria) si può usare in più programmi;
  \item \textbf{lavoro di gruppo}: persone diverse lavorano su file
        diversi;
  \item \textbf{compilazione più rapida}: se si modifica un solo file, si
        ricompila solo quello e si rifà il collegamento.
\end{itemize}

\subsection{Un piccolo progetto}

Prendiamo un programma che calcola l'area di un cerchio e di un quadrato, e
dividiamolo in tre file: le funzioni di calcolo in \code{aree.c}, il
programma principale in \code{main.c}, e un \textbf{file di intestazione}
\code{aree.h} con ciò che serve a entrambi.

\begin{lstlisting}[style=h, name=aree.h]
/* aree.h: costanti e prototipi delle funzioni sulle aree */

#define PI 3.14159

double area_cerchio(double raggio);
double area_quadrato(double lato);
\end{lstlisting}

\begin{affianco}[0.46]
\begin{lstlisting}[style=c, name=aree.c]
/* aree.c: definizioni delle funzioni sulle aree */

#include "aree.h"

double area_cerchio(double raggio)
{
    return PI * raggio * raggio;
}

double area_quadrato(double lato)
{
    return lato * lato;
}
\end{lstlisting}
\risultato
\begin{lstlisting}[style=c, name=main.c]
#include <stdio.h>
#include "aree.h"

int main(void)
{
    double x;

    printf("Inserisci una misura: ");
    scanf("%lf", &x);
    printf("cerchio di raggio %g: area %g\n",
           x, area_cerchio(x));
    printf("quadrato di lato %g: area %g\n",
           x, area_quadrato(x));
    return 0;
}
\end{lstlisting}
\end{affianco}

Il ruolo di ciascun file:

\begin{itemize}
  \item \code{aree.h} contiene le \textbf{dichiarazioni}: la costante e i
        prototipi. Chi lo include sa \textit{che} le funzioni esistono e come
        si chiamano;
  \item \code{aree.c} contiene le \textbf{definizioni}: il codice delle
        funzioni, scritto una volta sola. Include anche lui \code{aree.h},
        così il compilatore controlla che le definizioni corrispondano ai
        prototipi;
  \item \code{main.c} include \code{aree.h} e può quindi chiamare le
        funzioni, anche se il loro codice sta in un altro file.
\end{itemize}

\subsection{Compilare separatamente, collegare insieme}

Ogni file \code{.c} si compila per conto suo con l'opzione \code{-c}, che
produce il file oggetto; poi il linker unisce i file oggetto
nell'eseguibile:

\begin{console}
$ gcc -Wall -c main.c
$ gcc -Wall -c aree.c
$ gcc -o aree main.o aree.o
$ ./aree
Inserisci una misura: <<2>>
cerchio di raggio 2: area 12.5664
quadrato di lato 2: area 4
\end{console}

Che cosa contengono i due file oggetto, e come fa il linker a metterli
insieme? Lo mostra \code{nm} (capitolo~\ref{chap:dal-problema}): in
\code{main.o} le due funzioni sono \textit{undefined} (\code{U}), cioè usate
ma senza codice; in \code{aree.o} sono presenti (\code{T}). Il linker
abbina ogni nome mancante con la sua definizione (schema a destra).

\begin{affianco}[0.5][c]
\begin{console}
$ nm main.o
                 U __isoc99_scanf
                 U __stack_chk_fail
                 U area_cerchio
                 U area_quadrato
0000000000000000 T main
                 U printf
$ nm aree.o
0000000000000000 T area_cerchio
0000000000000025 T area_quadrato
\end{console}
\risultato
\begin{immagine}
\begin{tikzpicture}[font=\FiraCodeNL\scriptsize,
    sim/.style={anchor=west, inner sep=1.5pt},
    box/.style={draw=primary!65!black, fill=primary!4!white, line width=0.7pt,
      rounded corners=3pt}]
  \draw[box] (0,0) rectangle (2.9,-2.3);
  \node[font=\Lato\scriptsize\bfseries, anchor=south west] at (0,0.02) {main.o};
  \node[sim] at (0.1,-0.35) {T main};
  \node[sim, text=merr] (u1) at (0.1,-0.85) {U area\_cerchio};
  \node[sim, text=merr] (u2) at (0.1,-1.35) {U area\_quadrato};
  \node[sim, text=neutral!45!black] at (0.1,-1.85) {U printf, scanf\dots};
  \draw[box] (4.6,-0.4) rectangle (7.6,-1.8);
  \node[font=\Lato\scriptsize\bfseries, anchor=south west] at (4.6,-0.38) {aree.o};
  \node[sim, text=mok] (t1) at (4.7,-0.85) {T area\_cerchio};
  \node[sim, text=mok] (t2) at (4.7,-1.35) {T area\_quadrato};
  \draw[wtarrowacc] (u1.east) -- (t1.west);
  \draw[wtarrowacc] (u2.east) -- (t2.west);
  \node[mnota, anchor=north, align=center] at (3.8,-2.4) {il linker abbina ogni nome usato (\code{U})\\alla sua definizione (\code{T}); \code{printf}\\e \code{scanf} le trova nella libreria standard};
\end{tikzpicture}
\end{immagine}
\end{affianco}

Che cosa succede se si dimentica \code{aree.o} al momento del linking? La
compilazione di \code{main.c} era già riuscita (grazie a \code{aree.h} il
compilatore sapeva tutto ciò che serviva), ma il linker non trova il codice
delle due funzioni:

\begin{console}
$ gcc -o aree main.o
/usr/bin/ld: main.o: in function `main':
main.c:(.text+0x54): undefined reference to `area_cerchio'
/usr/bin/ld: main.c:(.text+0x89): undefined reference to `area_quadrato'
collect2: error: ld returned 1 exit status
\end{console}

È lo stesso errore di \code{pow} senza \code{-lm} del
capitolo~\ref{chap:primi-programmi}: \textit{undefined reference} significa
sempre “il linker non trova il codice di questa funzione”.

\section{Il progetto delle lezioni}

\subsection{I file}

A lezione si usa un progetto su tre file, che gestisce un elenco di esami
sostenuti da alcuni studenti (letti da un file di testo): \code{main.c},
\code{funzioni.c} e \code{studente.h} (figura~\ref{fig:progetto}). Il codice
usa strumenti che vedremo più avanti (strutture, liste, file), e verrà
spiegato per intero nel capitolo~\ref{chap:liste}; per capire \code{make}
basta sapere chi dipende da chi.

\begin{figure}[H]
  \centering
  \begin{tikzpicture}[font=\Lato\scriptsize,
      f/.style={wtnodeplain, minimum width=3.2cm, minimum height=1.15cm, align=center}]
    \node[f, draw=accent!60!black, fill=accent!7!white] (h) at (0,0) {{\FiraCodeNL\bfseries studente.h}\\la struttura \code{studente}\\e i prototipi delle funzioni};
    \node[f] (m) at (-3.6,-1.9) {{\FiraCodeNL\bfseries main.c}\\legge i dati, filtra,\\scrive il risultato};
    \node[f] (fu) at (3.6,-1.9) {{\FiraCodeNL\bfseries funzioni.c}\\le funzioni sulla lista\\di studenti};
    \draw[wtarrowacc] (h) -- node[wtlab, above left, fill=none] {incluso da} (m);
    \draw[wtarrowacc] (h) -- node[wtlab, above right, fill=none] {incluso da} (fu);
    \draw[wtarrow, dashed] (m) -- node[wtlab, below=2pt, fill=none] {chiama le funzioni di} (fu);
  \end{tikzpicture}
  \caption{Il progetto delle lezioni: due sorgenti e un header condiviso.}
  \label{fig:progetto}
\end{figure}

L'eseguibile si crea come abbiamo visto: la compilazione
(\code{gcc -c main.c funzioni.c}, che si può fare anche con un solo comando
per più file) produce \code{main.o} e \code{funzioni.o}; il linking
(\code{gcc -o esercizio main.o funzioni.o}) produce l'eseguibile
\code{esercizio}.

\subsection{Il problema}

Con molti file la compilazione può essere lunga, e non sempre è necessario
ricompilare tutti i sorgenti: se si modifica solo \code{funzioni.c},
\code{main.o} è ancora valido. Vorremmo rendere le operazioni più semplici
e automatiche.

Un primo passo è uno \term{script}: un file di testo che contiene i comandi,
da eseguire tutti in fila.

\begin{affianco}[0.5][c]
\begin{lstlisting}[style=shell, name=compila.sh]
#!/bin/bash
gcc -c main.c
gcc -c funzioni.c
gcc -o esercizio main.o funzioni.o
\end{lstlisting}
\risultato
\begin{console}
$ chmod +x compila.sh
$ ./compila.sh
\end{console}
\end{affianco}

La prima riga (\code{\#!/bin/bash}, detta \textit{shebang}) dice al sistema
quale programma deve eseguire lo script: la shell bash. \code{chmod +x}
rende il file eseguibile, una volta per tutte; da lì in poi
\code{./compila.sh} esegue i tre comandi. Il procedimento è automatico, ma i
sorgenti vengono comunque compilati \textbf{tutti}, ogni volta. Vorremmo un
meccanismo che compila solo i file che sono stati modificati.

\section{make}

\subsection{L'idea: dipendenze e date}

Il comando \code{make} esegue la compilazione in modo automatico seguendo,
fondamentalmente, due regole:

\begin{itemize}
  \item la \textbf{dipendenza} fra i file: \code{esercizio} si costruisce da
        \code{main.o} e \code{funzioni.o}; \code{main.o} si costruisce da
        \code{main.c} (e da \code{studente.h}, che \code{main.c} include), e
        così via (figura~\ref{fig:dipendenze});
  \item l'\textbf{ordine temporale} in cui sono stati creati o modificati i
        file sorgente, oggetto ed eseguibile: ogni file ha una data e ora di
        ultima modifica, che \code{ls -l} mostra.
\end{itemize}

\begin{figure}[H]
  \centering
  \begin{tikzpicture}[font=\Lato\scriptsize,
      t/.style={wtnode, font=\FiraCodeNL\scriptsize, minimum width=2.3cm},
      s/.style={wtnodeplain, font=\FiraCodeNL\scriptsize, minimum width=2.3cm},
      ora/.style={font=\Lato\tiny, text=neutral!40!black}]
    \node[t, draw=red-gradient-6!70!black, fill=red-gradient-1!20!white] (e) at (0,0) {esercizio};
    \node[t] (mo) at (-3.0,-1.5) {main.o};
    \node[t] (fo) at (3.0,-1.5) {funzioni.o};
    \node[s] (mc) at (-4.6,-3.0) {main.c};
    \node[s, draw=accent!60!black, fill=accent!7!white] (h) at (0,-3.0) {studente.h};
    \node[s] (fc) at (4.6,-3.0) {funzioni.c};
    \foreach \a/\b in {e/mo, e/fo, mo/mc, mo/h, fo/h, fo/fc}
      \draw[wtarrow, <-] (\a) -- (\b);
    \node[ora, anchor=west] at ([xshift=3pt]e.east) {modificato alle 10:07};
    \node[ora, anchor=east] at ([xshift=-3pt]mo.west) {10:05};
    \node[ora, anchor=west] at ([xshift=3pt]fo.east) {10:06};
    \node[ora, anchor=north] at ([yshift=-2pt]mc.south) {10:00};
    \node[ora, anchor=north] at ([yshift=-2pt]h.south) {9:55};
    \node[ora, anchor=north, text=merr] at ([yshift=-2pt]fc.south) {\textbf{10:12}: modificato dopo \code{funzioni.o}!};
  \end{tikzpicture}
  \caption{Il grafo delle dipendenze del progetto: le frecce vanno da ciò
           che serve a ciò che si costruisce. \code{funzioni.c} è più
           recente di \code{funzioni.o}, quindi \code{funzioni.o} va
           ricostruito, e di conseguenza anche \code{esercizio}.}
  \label{fig:dipendenze}
\end{figure}

Il ragionamento di \code{make} è quello che faremmo noi: per costruire un
file, prima si assicura che siano aggiornati i file da cui dipende (e così
via, all'indietro); poi, se il file \textbf{non esiste} oppure è \textbf{più
vecchio} di almeno uno dei file da cui dipende, lo ricostruisce eseguendo i
comandi indicati. Altrimenti non fa niente. Nella situazione della
figura~\ref{fig:dipendenze}, \code{make} ricompila solo \code{funzioni.c} e
rifà il linking; \code{main.c} non viene toccato.

\subsection{Il makefile}

\code{make} esegue le operazioni descritte in un file, chiamato di solito
\code{makefile} o \code{Makefile}, che contiene la descrizione delle
dipendenze, i comandi da eseguire e, eventualmente, delle definizioni di
\textit{macro}. Per usarlo basta scrivere \code{make} nella cartella che lo
contiene; se il file ha un altro nome, per esempio \code{makeEsercizio}, lo
si indica con l'opzione \code{-f}:

\begin{console}
$ make
$ make -f makeEsercizio
\end{console}

Ecco il primo makefile delle lezioni. Nei listati di questo capitolo il
carattere di tabulazione (TAB) è disegnato come una freccia
\tikz[baseline=-2.6pt]\draw[vsclinenum, -{Stealth[length=3pt,width=3pt]}, line width=0.5pt] (0,0) -- (0.5,0);,
come fa VS Code quando si chiede di mostrare gli spazi bianchi: il motivo
sarà chiaro fra poco.

\begin{lstlisting}[style=make, name=makefile1]
# makefile di esempio
# versione 1
#

eseguibile: main.o funzioni.o
(*@\TAB@*)gcc -o esercizio main.o funzioni.o

main.o: main.c
(*@\TAB@*)gcc -c main.c

funzioni.o: funzioni.c
(*@\TAB@*)gcc -c funzioni.c
\end{lstlisting}

\subsection{I tipi di riga}

In un makefile si distinguono quattro tipi di righe:

\begin{itemize}
  \item le righe di \textbf{dipendenza}: definiscono una relazione fra un
        \term{target} (l'obiettivo, ciò che si costruisce) e uno o più file
        da cui dipende, separati dal carattere \code{:}. Per esempio
        \code{main.o: main.c};
  \item le righe di \textbf{comando}: definiscono l'operazione da eseguire
        per ottenere il target, e seguono la sua riga di dipendenza.
        \textbf{Devono cominciare con un TAB};
  \item le righe di \textbf{commento}: cominciano con \code{\#};
  \item le righe di definizione di \textbf{macro}, che vedremo fra poco.
\end{itemize}

\error{Il TAB non è uno spazio}{
  Le righe di comando devono cominciare con un carattere di
  \textbf{tabulazione}, non con degli spazi, anche se a occhio sono
  identici. Il file \code{makefile4} delle lezioni contiene proprio questo
  errore (quattro spazi al posto del TAB), e \code{make} si rifiuta di
  proseguire con un messaggio poco chiaro:
  \textit{makefile4:9: *** missing separator. Stop.} VS Code inserisce
  davvero un TAB solo se riconosce il file come makefile, cioè se il file si
  chiama \code{makefile} o \code{Makefile}: in basso a destra, nella barra
  di stato, compare allora \textit{Tab Size} invece di \textit{Spaces}.}

\subsection{Che cosa fa \code{make}, passo per passo}

Proviamo il primo makefile. La prima volta mancano tutti i file oggetto e
l'eseguibile, quindi \code{make} esegue tutto, cominciando dai file da cui
dipende il primo target:

\begin{console}
$ make -f makefile1
gcc -c main.c
gcc -c funzioni.c
gcc -o esercizio main.o funzioni.o
\end{console}

\code{make} stampa ogni comando prima di eseguirlo. Si noti l'ordine: per
costruire \code{eseguibile} servono \code{main.o} e \code{funzioni.o}, quindi
prima vengono costruiti loro. Se ora modifichiamo soltanto
\code{funzioni.c} (il comando \code{touch} cambia la data di un file, come
se lo avessimo modificato e salvato) e rilanciamo \code{make}, viene
ricompilato solo quello:

\begin{console}
$ touch funzioni.c
$ make -f makefile1
gcc -c funzioni.c
gcc -o esercizio main.o funzioni.o
\end{console}

\subsection{Le dipendenze dall'header}

Che cosa succede se modifichiamo \code{studente.h}? Con il primo makefile,
niente di buono:

\begin{console}
$ touch studente.h
$ make -f makefile1
gcc -o esercizio main.o funzioni.o
\end{console}

Nessun file viene ricompilato, perché nel makefile i file oggetto dipendono
solo dai rispettivi \code{.c}. Ma entrambi i \code{.c} \textbf{includono}
\code{studente.h}: se nell'header cambia, per esempio, la struttura
\code{studente}, i due file oggetto sono stati compilati con la versione
vecchia e il programma risultante è incoerente, con errori imprevedibili.
La soluzione è dichiarare la dipendenza:

\begin{lstlisting}[style=make, name=makefile2]
# makefile di esempio
# versione 2

eseguibile: main.o funzioni.o
(*@\TAB@*)gcc -o esercizio main.o funzioni.o

main.o: main.c studente.h
(*@\TAB@*)gcc -c main.c

funzioni.o: funzioni.c studente.h
(*@\TAB@*)gcc -c funzioni.c
\end{lstlisting}

\begin{console}
$ make -f makefile2
gcc -c main.c
gcc -c funzioni.c
gcc -o esercizio main.o funzioni.o
\end{console}

Ora, dopo una modifica all'header, vengono ricompilati entrambi i sorgenti.
La regola generale: \textbf{un file oggetto dipende dal suo \code{.c} e da
tutti gli header (nostri) che questo include}.

\subsection{Il nome del target}

Un'altra stranezza: con i primi due makefile il linking viene eseguito
\textbf{ogni volta}, anche se nulla è cambiato. Possiamo evitarlo? Il motivo
è il nome del target: \code{eseguibile}. \code{make} controlla se esiste un
file con quel nome e se è aggiornato; ma il comando produce un file chiamato
\code{esercizio}, quindi un file \code{eseguibile} non esiste mai, e
\code{make} rifà il linking per crearlo. Basta che il target abbia lo stesso
nome del file prodotto:

\begin{lstlisting}[style=make, name=makefile3]
# makefile di esempio
# versione 3

esercizio: main.o funzioni.o
(*@\TAB@*)gcc -o esercizio main.o funzioni.o

main.o: main.c studente.h
(*@\TAB@*)gcc -c main.c

funzioni.o: funzioni.c studente.h
(*@\TAB@*)gcc -c funzioni.c
\end{lstlisting}

\begin{console}
$ make -f makefile3
make: 'esercizio' is up to date.
\end{console}

Tutto è aggiornato, e \code{make} lo dice senza eseguire nulla.

\subsection{Le macro}

Una riga di definizione di \term{macro} definisce un nome simbolico per un
insieme di oggetti, con la sintassi \code{NOME = valore}; il nome si usa poi
scrivendo \code{\$(NOME)}. È l'equivalente del \code{\#define} del C: si
scrive l'elenco una volta sola, e se cambia si modifica in un solo punto.

\begin{lstlisting}[style=make, name=makefile4]
# makefile di esempio
# versione 4

OBJ = funzioni.o main.o

esercizio: $(OBJ)
(*@\TAB@*)gcc -o esercizio $(OBJ)

main.o: main.c studente.h
(*@\TAB@*)gcc -c main.c

funzioni.o: funzioni.c studente.h
(*@\TAB@*)gcc -c funzioni.c
\end{lstlisting}

(Questa è la versione corretta, con il TAB; nel file distribuito con le
lezioni la riga del linking comincia con quattro spazi e produce l'errore
\textit{missing separator}.)

\section{Un makefile completo}

Nella pratica i makefile usano quasi sempre qualche macro in più e un
target di servizio per fare pulizia. Ecco un makefile per il progetto
\code{aree} di inizio capitolo:

\begin{affianco}[0.5]
\begin{lstlisting}[style=make, name=Makefile]
CC = gcc
CFLAGS = -Wall
OBJ = main.o aree.o

aree: $(OBJ)
(*@\TAB@*)$(CC) $(CFLAGS) -o aree $(OBJ)

main.o: main.c aree.h
(*@\TAB@*)$(CC) $(CFLAGS) -c main.c

aree.o: aree.c aree.h
(*@\TAB@*)$(CC) $(CFLAGS) -c aree.c

clean:
(*@\TAB@*)rm -f aree $(OBJ)
\end{lstlisting}
\risultato
\begin{console}
$ make
gcc -Wall -c main.c
gcc -Wall -c aree.c
gcc -Wall -o aree main.o aree.o
$ make
make: 'aree' is up to date.
$ touch main.c
$ make
gcc -Wall -c main.c
gcc -Wall -o aree main.o aree.o
$ make clean
rm -f aree main.o aree.o
\end{console}
\end{affianco}

\begin{itemize}
  \item \code{CC} e \code{CFLAGS} sono i nomi tradizionali delle macro per
        il compilatore e le sue opzioni: per aggiungere un'opzione a tutte
        le compilazioni (per esempio \code{-g}, che serve al debugger) basta
        cambiare una riga;
  \item senza argomenti, \code{make} costruisce il \textbf{primo} target del
        file (qui \code{aree}); con \code{make main.o} si costruisce solo
        quello;
  \item \code{clean} non corrisponde a nessun file e non ha dipendenze: è un
        target “di servizio”, che con \code{make clean} cancella
        l'eseguibile e i file oggetto, per ripartire da zero.
\end{itemize}

\info{make non è solo per il C}{
  \code{make} non sa nulla di C: sa solo che certi file si costruiscono da
  altri file eseguendo dei comandi, se sono più vecchi. Per questo si usa
  anche per compilare documenti, generare siti, elaborare dati: ovunque ci
  sia una catena di file da ricostruire quando qualcosa cambia.}

\insintesi{
  Un programma si divide in più file \code{.c}, ciascuno compilato
  separatamente in un file oggetto (\code{gcc -c}), e in header \code{.h}
  con le dichiarazioni condivise; il linker unisce gli oggetti, abbinando
  ogni funzione usata alla sua definizione (\textit{undefined reference}
  vuol dire che una definizione manca). \code{make} legge un
  \textbf{makefile} fatto di righe di dipendenza (\code{target: file}),
  righe di comando (che cominciano con un \textbf{TAB}), commenti e macro, e
  ricostruisce un target solo se non esiste o è più vecchio di qualcosa da
  cui dipende. Gli header vanno elencati fra le dipendenze, e il nome del
  target deve essere quello del file prodotto.}

\section{Riferimenti}

\begin{itemize}
  \item \textbf{GNU make --- An Introduction to Makefiles}, il manuale
        ufficiale. \\
        \urlx{https://www.gnu.org/software/make/manual/html_node/Introduction.html}
  \item \textbf{Deitel \& Deitel}, \textit{Il linguaggio C}: appendice sulla
        compilazione di programmi su più file.
\end{itemize}
